\documentclass{IEEEtran}
\usepackage[utf8]{vietnam}

\usepackage{cite}
\usepackage{amsmath,amssymb,amsfonts}
\usepackage{graphicx}
\usepackage{textcomp,nicefrac}
% --- Goi bo sung (khong anh huong dinh dang IEEEtran) ---
\usepackage{xcolor}
\usepackage{tikz}
\usetikzlibrary{positioning,arrows.meta,shapes.geometric}
\usepackage{algorithm}
\usepackage{algpseudocode}
% Danh dau cho can bo sung so lieu / thong tin (xoa truoc khi nop)
\newcommand{\todo}[1]{\textcolor{red}{[#1]}}

\def\BibTeX{{\rm B\kern-.05em{\sc i\kern-.025em b}\kern-.08em
T\kern-.1667em\lower.7ex\hbox{E}\kern-.125emX}}
\markboth{STUDENT SCIENTIFIC RESEARCH COMMUNICATION, Volume II, 2026}
{Nguyễn Hữu Sang \MakeLowercase{\textit{et al.}}: Pipeline lồng tiếng tự động cho video bài giảng với điều chỉnh tốc độ đọc thích ứng \newline (November 2026)}

\begin{document}
\title{Pipeline lồng tiếng tự động cho video bài giảng trực tuyến với cơ chế điều chỉnh tốc độ đọc thích ứng}

\author{Nguyễn Hữu Sang, Nguyễn Chí Thiện, và Nguyễn Thị Hoàng Khánh
\thanks{Bài báo được gửi ngày 14/11/2026 tới Hội nghị SSRC-II'2026.}
\thanks{Nguyễn Hữu Sang và Nguyễn Chí Thiện là sinh viên Khoa Công nghệ Thông tin, Trường Đại học Công nghiệp Thành phố Hồ Chí Minh (e-mail: 24huusang6a2@gmail.com, nguyenthien110120@gmail.com).}
\thanks{ThS. Nguyễn Thị Hoàng Khánh là giảng viên Khoa Công nghệ Thông tin, Trường Đại học Công nghiệp Thành phố Hồ Chí Minh (e-mail: ngthihoangkhanh@gmail.com).}}

\maketitle

\begin{abstract}
Các khóa học trực tuyến quốc tế phần lớn được giảng bằng tiếng nước ngoài, buộc người học Việt Nam vừa nghe vừa đọc phụ đề. Bài báo đề xuất một pipeline lồng tiếng tự động cho video bài giảng, tích hợp trong hệ thống quản lý học tập: nhận dạng giọng nói có mốc thời gian, dịch bằng mô hình ngôn ngữ lớn có giới hạn độ dài từng câu, tổng hợp giọng nói và ghép âm thanh theo vị trí của từng câu gốc. Đóng góp chính là thuật toán điều chỉnh tốc độ đọc thích ứng: hệ thống so thời lượng âm thanh tổng hợp của mỗi câu với khung thời gian gốc, rút gọn theo lô các câu vượt ngưỡng rồi tăng tốc tín hiệu tối đa ba mươi phần trăm; các câu được xếp theo mốc thời gian tuyệt đối để độ trễ không cộng dồn. Video được chia thành các đoạn mười phút xử lý độc lập. Thực nghiệm trên bốn bài giảng công nghệ thông tin, năm mươi phút, cho thấy độ lệch thời lượng trung bình mỗi câu giảm từ 1,89 giây (dịch không ràng buộc) xuống 0,19 giây; âm thanh sau khi ghép có thời lượng đúng bằng video gốc với độ trễ tuyệt đối trung bình 0,40 giây. Trên hệ thống thật, mỗi phút video được xử lý trong khoảng mười sáu giây. Nghiên cứu chưa đánh giá cảm nhận của người nghe.
\end{abstract}

\begin{IEEEkeywords}
Điều chỉnh tốc độ đọc, hệ thống quản lý học tập, lồng tiếng tự động, mô hình ngôn ngữ lớn, nhận dạng giọng nói, tổng hợp giọng nói.
\end{IEEEkeywords}

% =====================================================================
\section{Giới thiệu}
\label{sec:introduction}
\IEEEPARstart{H}{ọc} trực tuyến đã trở thành kênh tiếp cận tri thức phổ biến, trong đó phần lớn khóa học chất lượng cao trên các nền tảng quốc tế được giảng bằng tiếng Anh. Với người học Việt Nam, rào cản ngôn ngữ thường được khắc phục bằng phụ đề. Tuy nhiên, theo nguyên lý dư thừa và chia tách sự chú ý trong học tập đa phương tiện \cite{mayer2009,kalyuga1999}, việc buộc người học vừa đọc chữ vừa quan sát hình minh họa, mã nguồn hay thao tác trên màn hình làm tăng tải nhận thức và giảm khả năng ghi nhớ. Lồng tiếng sang ngôn ngữ mẹ đẻ giúp người học dồn toàn bộ sự chú ý thị giác vào nội dung bài giảng.

Lồng tiếng thủ công tốn kém và không thể áp dụng cho hàng nghìn video bài giảng. Sự phát triển của các mô hình nhận dạng giọng nói quy mô lớn \cite{radford2023whisper}, mô hình ngôn ngữ lớn (LLM) \cite{gemini2023} và tổng hợp giọng nói (TTS) thần kinh mở ra khả năng tự động hóa toàn bộ quy trình. Dù vậy, lồng tiếng tự động khác dịch phụ đề ở một ràng buộc quan trọng: câu nói được lồng tiếng phải \emph{khớp thời lượng} với câu gốc \cite{federico2020,brannon2023}. Bản dịch tiếng Việt thường dài hơn câu tiếng Anh tương ứng; nếu phát nguyên văn, âm thanh sẽ trễ dần so với hình ảnh, còn nếu tăng tốc quá mức thì giọng đọc mất tự nhiên.

Bài báo này trình bày một pipeline lồng tiếng tự động cho video bài giảng, được xây dựng và tích hợp trong hệ thống quản lý học tập (LMS) LinguaLearn. Các đóng góp chính gồm:
\begin{itemize}
\item Thuật toán \emph{điều chỉnh tốc độ đọc thích ứng} (Adaptive Speech Rate) kết hợp ba chiến lược theo từng câu: để khoảng lặng, rút gọn theo lô bằng LLM có giới hạn số vòng, và tăng tốc tín hiệu bằng FFmpeg \texttt{atempo} trong giới hạn 30\%.
\item Quy trình \emph{dịch có ràng buộc độ dài} trong một lời gọi LLM cho mỗi đoạn: mỗi câu được gửi kèm ngân sách thời lượng và giới hạn số ký tự ước theo tốc độ đọc đo được của giọng tổng hợp, không cần huấn luyện lại mô hình dịch.
\item Kiến trúc xử lý \emph{phân đoạn theo lớp} (Layered Chunking) mười phút, cho phép báo tiến độ thời gian thực, chịu lỗi cục bộ, cùng các cơ chế chống trùng tác vụ và tái sử dụng kết quả nhằm giảm chi phí gọi dịch vụ AI.
\item Đánh giá thực nghiệm về độ khớp thời lượng, thời gian xử lý và chi phí trên tập video bài giảng thực tế.
\end{itemize}

Phần còn lại của bài báo được tổ chức như sau: Mục~\ref{sec:related} trình bày các nghiên cứu liên quan; Mục~\ref{sec:method} mô tả pipeline và thuật toán đề xuất; Mục~\ref{sec:impl} trình bày cách triển khai trong hệ thống; Mục~\ref{sec:exp} trình bày thực nghiệm và Mục~\ref{sec:conclusion} kết luận.

% =====================================================================
\section{Nghiên cứu liên quan}
\label{sec:related}

\subsection{Lồng tiếng tự động}
Federico và cộng sự \cite{federico2020} mở rộng bài toán dịch tiếng nói sang tiếng nói thành lồng tiếng tự động, trong đó yêu cầu đồng bộ nhịp (isochrony) giữa câu gốc và câu dịch được đặt lên hàng đầu. \"{O}ktem và cộng sự \cite{oktem2019} căn chỉnh các cụm ngữ điệu giữa hai ngôn ngữ để phân bổ lời thoại khớp với các khoảng ngắt. Nghiên cứu thực tiễn của Brannon và cộng sự \cite{brannon2023} trên các bản lồng tiếng chuyên nghiệp chỉ ra rằng người lồng tiếng ưu tiên độ trung thành nội dung và sự tự nhiên hơn là khớp thời lượng tuyệt đối, đồng thời thường xuyên điều chỉnh tốc độ nói.

\subsection{Điều khiển độ dài bản dịch}
Một hướng tiếp cận là điều khiển độ dài ngay trong mô hình dịch. Lakew và cộng sự \cite{lakew2022} đề xuất dịch máy đẳng cự (isometric MT) để bản dịch có số ký tự tương đương câu nguồn. Wu và cộng sự \cite{wu2023videodubber} đưa thông tin thời lượng tiếng nói vào mô hình dịch để kiểm soát độ dài ở mức âm vị. Các phương pháp này đòi hỏi huấn luyện lại mô hình dịch, khó áp dụng khi hệ thống sử dụng LLM thông qua API.

\subsection{Nhận dạng giọng nói và giải pháp thương mại}
Whisper \cite{radford2023whisper} cho độ chính xác cao trên nhiều ngôn ngữ; WhisperX \cite{bain2023whisperx} bổ sung căn chỉnh mốc thời gian cấp độ từ cho âm thanh dài. Về sản phẩm, một số nền tảng chia sẻ video và dịch vụ thương mại đã cung cấp tính năng lồng tiếng tự động, tuy nhiên chúng tách rời khỏi hệ thống quản lý học tập, tính phí theo phút video và không công bố chi tiết thuật toán đồng bộ thời lượng.

Khác với các nghiên cứu trên, bài báo không huấn luyện mô hình mới mà tập trung vào \emph{điều phối} các dịch vụ AI có sẵn: ràng buộc thời lượng được xử lý ở hai tầng là tầng prompt (ngân sách thời lượng và giới hạn số ký tự, gần với ý tưởng dịch đẳng cự \cite{lakew2022} nhưng thực hiện qua lời nhắc thay vì huấn luyện) và tầng hậu xử lý (rút gọn có điều kiện và điều chỉnh tốc độ đọc).

% =====================================================================
\section{Phương pháp đề xuất}
\label{sec:method}

\subsection{Tổng quan pipeline}
Hình~\ref{fig:pipeline} mô tả pipeline đề xuất. Đầu vào là một video bài giảng (tệp tải lên hoặc đường dẫn YouTube) và ngôn ngữ đích do học viên chọn. Đầu ra là một tệp âm thanh lồng tiếng duy nhất có cùng thời lượng với video, được phát đồng bộ với video gốc đã tắt tiếng.

\begin{figure}[t]
\centering
\begin{tikzpicture}[
  node distance=3.2mm,
  box/.style={draw, rounded corners=2pt, text width=7.0cm, minimum height=6mm, align=center, font=\footnotesize},
  dec/.style={draw, diamond, aspect=2.6, inner sep=1pt, align=center, font=\scriptsize},
  arr/.style={-{Stealth[length=2mm]}}]
\node[box] (in) {Video bài giảng + ngôn ngữ đích};
\node[box, below=of in] (chunk) {Tách âm thanh, chia đoạn cố định 10 phút};
\node[box, below=of chunk] (asr) {Nhận dạng giọng nói (Whisper): câu + mốc thời gian\\ \scriptsize (bỏ qua nếu đã có bản ký âm từ tác vụ trước)};
\node[box, below=of asr] (tr) {Dịch bằng LLM (một lời gọi mỗi đoạn):\\ ngân sách thời lượng + giới hạn số ký tự};
\node[box, below=of tr] (tts) {Tổng hợp giọng nói song song (Azure TTS)};
\node[box, below=of tts] (asr2) {Điều chỉnh tốc độ đọc thích ứng (Thuật toán~\ref{alg:asr})};
\node[box, below=of asr2] (cat) {Xếp câu theo mốc tuyệt đối $\rightarrow$ âm thanh của đoạn};
\node[box, below=of cat] (out) {Ghép các đoạn $\rightarrow$ tệp lồng tiếng cuối cùng};
\foreach \a/\b in {in/chunk,chunk/asr,asr/tr,tr/tts,tts/asr2,asr2/cat,cat/out}
  \draw[arr] (\a) -- (\b);
\draw[arr, dashed] (cat.east) -- ++(4mm,0) |- node[pos=0.25, right, font=\scriptsize, align=left]{lặp\\từng\\đoạn} (asr.east);
\end{tikzpicture}
\caption{Pipeline lồng tiếng tự động đề xuất. Mỗi đoạn 10 phút được xử lý độc lập và báo tiến độ ngay khi hoàn thành.}
\label{fig:pipeline}
\end{figure}

\subsection{Phân đoạn theo lớp}
Video dài (tối đa 180 phút) nếu xử lý nguyên khối sẽ khiến người học phải chờ đến khi toàn bộ pipeline kết thúc, đồng thời một lỗi nhỏ có thể làm hỏng cả tác vụ. Hệ thống chia âm thanh thành $K = \lceil D / L \rceil$ đoạn cố định, với $D$ là thời lượng video và $L = 600$ giây; đoạn cuối chứa phần dư. Mỗi đoạn được xử lý độc lập qua các bước nhận dạng, dịch và tổng hợp; khi một đoạn hoàn thành, tiến độ được phát qua kênh thời gian thực để giao diện cập nhật ngay. Với bài có nhiều đoạn, các đoạn được xử lý theo \emph{dây chuyền hai công đoạn}: trong khi đoạn $k$ đang tổng hợp giọng nói và tải lên, đoạn $k+1$ đã được nhận dạng và dịch; kết quả vẫn được báo về theo đúng thứ tự đoạn. Cách này rút ngắn tổng thời gian với bài dài hơn một đoạn nhưng không đổi thời gian tới khi có đoạn đầu tiên.

Mỗi đoạn được thử lại tối đa ba lần với thời gian chờ tăng theo hàm mũ. Nếu vẫn thất bại, hệ thống \emph{giữ âm thanh gốc} cho riêng đoạn đó thay vì hủy toàn bộ tác vụ; người học vẫn nghe được bản lồng tiếng ở các đoạn còn lại. Đoạn có tỷ lệ lời thoại dưới 10\% thời lượng (ví dụ phần trình chiếu không lời) được bỏ qua bước dịch và giữ nguyên âm thanh gốc.

\subsection{Nhận dạng giọng nói}
Âm thanh của từng đoạn được gửi tới mô hình Whisper \cite{radford2023whisper} (biến thể large-v3-turbo) để nhận về danh sách câu $s_i$ kèm mốc kết thúc $e_i$, mốc thời gian theo từng từ, và ngôn ngữ nguồn được tự động phát hiện. Mốc bắt đầu $b_i$ được lấy theo \emph{từ đầu tiên} của câu thay vì mốc đầu đoạn của Whisper, vì mốc cấp câu thường gộp cả khoảng không lời phía trước (ví dụ nhạc mở đầu: câu đầu của một bài Crash Course được ghi bắt đầu ở 0,0 giây trong khi người nói bắt đầu ở 3,16 giây), khiến lời lồng tiếng đè lên phần không lời. Khung thời gian của câu là $T_i^{\mathrm{orig}} = e_i - b_i$. Câu chỉ gồm dấu câu (thường gặp ở đoạn nhạc) không được gửi tới TTS mà thay bằng khoảng lặng ngắn. Bản ký âm gốc được lưu vào cơ sở dữ liệu; khi bài giảng được lồng tiếng sang ngôn ngữ đích khác, bước nhận dạng được bỏ qua hoàn toàn.

\subsection{Dịch có ràng buộc độ dài}
\label{sub:translate}
Thay vì dịch từng câu riêng lẻ, toàn bộ các câu trong một đoạn được đánh số và gửi trong \emph{một} lời gọi LLM để mô hình nắm được ngữ cảnh bài giảng, giúp thuật ngữ được dịch nhất quán. Lời nhắc yêu cầu dịch đúng nghĩa, giữ thuật ngữ chuyên môn và viết ngay thành lời nói tự nhiên khi đọc bằng TTS; đầu ra ở định dạng JSON giữ nguyên số dòng để kiểm tra tự động. Mỗi câu được gửi kèm ngân sách thời lượng $T_i^{\mathrm{orig}}$ và, khi có, giới hạn số ký tự
\begin{equation}
C_i = \bigl\lfloor \hat{c}_v \cdot T_i^{\mathrm{orig}} \cdot \rho \bigr\rfloor, \qquad \rho = 1{,}15,
\label{eq:maxchars}
\end{equation}
trong đó $\hat{c}_v$ là tốc độ đọc (ký tự trên giây) của giọng tổng hợp $v$. Tốc độ này không cố định theo ngôn ngữ mà được \emph{đo} từ lượt tổng hợp đầu của mỗi tác vụ và cập nhật trơn: $\hat{c}_v \gets 0{,}7\,\hat{c}_v + 0{,}3\,c_v^{\mathrm{đo}}$. Hệ số $\rho$ nhắm bản dịch tới $R \approx 1{,}15$, để phần dư nhỏ được xử lý bằng tăng tốc thay vì rút gọn. Ở tác vụ đầu tiên dùng một giọng, chưa có $\hat{c}_v$ nên bản dịch chỉ kèm ngân sách thời lượng. Câu bị LLM bỏ sót được dịch lại riêng; nếu vẫn thiếu thì giữ câu gốc để không có câu rỗng. Tên mô hình LLM và mức suy nghĩ được cấu hình qua biến môi trường, kèm danh sách mô hình dự phòng khi mô hình chính bị giới hạn tần suất gọi.

\subsection{Điều chỉnh tốc độ đọc thích ứng}
Các câu của một đoạn được tổng hợp giọng nói \emph{song song} ở tốc độ mặc định (tối đa 16 yêu cầu đồng thời). Với mỗi câu dịch $t_i$, ta đo thời lượng thực tế $T_i^{\mathrm{tts}}$ và tính tỷ số tốc độ:
\begin{equation}
R_i = \frac{T_i^{\mathrm{tts}}}{T_i^{\mathrm{orig}}}.
\label{eq:ratio}
\end{equation}
Với hai ngưỡng $R_{\min} = 1{,}0$ và $R_{\max} = 1{,}3$, hành động cho câu $i$ được chọn như sau:
\begin{equation}
A(R_i) =
\begin{cases}
\text{giữ nguyên}, & R_i \le R_{\min}\\
\text{tăng tốc } \times R_i, & R_{\min} < R_i \le R_{\max}\\
\text{rút gọn, rồi tăng tốc}, & R_i > R_{\max}
\end{cases}
\label{eq:policy}
\end{equation}

Tăng tốc được thực hiện trực tiếp trên tệp âm thanh đã tổng hợp bằng bộ lọc \texttt{atempo} của FFmpeg với hệ số $\min(R_i, R_{\max})$, giữ nguyên cao độ. Cách này không phải gọi lại dịch vụ TTS chỉ để đổi tốc độ đọc, nhưng là kéo giãn tín hiệu nên có thể ảnh hưởng tới độ tự nhiên khác với việc yêu cầu TTS đọc nhanh hơn. Ngưỡng 30\% được chọn theo kinh nghiệm vì vượt mức này giọng đọc dễ trở nên khó nghe; nghiên cứu này chưa kiểm chứng ngưỡng bằng đánh giá của người nghe.

Ở nhánh thứ ba, các câu vượt ngưỡng của cả đoạn được gửi \emph{theo lô} trong một lời gọi LLM, mỗi câu kèm giới hạn số ký tự $\lfloor |t_i| \cdot 1{,}25 \cdot T_i^{\mathrm{orig}} / T_i^{\mathrm{tts}} \rfloor$ ước theo tốc độ đọc của chính câu đó, nhắm $R \approx 1{,}25$ để chừa biên dưới trần. Các câu đã rút gọn được tổng hợp lại song song và tính lại $R_i$. Số vòng rút gọn bị giới hạn bởi $M = 2$ để tránh vòng lặp vô hạn; câu vẫn vượt ngưỡng sau $M$ vòng (hoặc khi lời gọi rút gọn lỗi) được tăng tốc tối đa $\times 1{,}3$ và ghi cảnh báo để quản trị viên xem xét. Câu có khung gốc rất ngắn ($T_i^{\mathrm{orig}} < T_{\min} = 2$ giây, ví dụ lời chào hay câu cảm ơn) không được rút gọn: thực nghiệm sơ bộ cho thấy phần lớn các câu này rút gọn hai lần vẫn vượt ngưỡng, chỉ làm tăng số lần gọi LLM. Câu ngắn vẫn được lồng tiếng ở tốc độ tối đa và phần dư được bù bằng khoảng lặng phía sau nhờ cơ chế xếp câu theo mốc tuyệt đối (Mục~\ref{sub:concat}). Thuật toán~\ref{alg:asr} tóm tắt quy trình cho một đoạn.

\begin{algorithm}[t]
\caption{Điều chỉnh tốc độ đọc thích ứng cho các câu của một đoạn}
\label{alg:asr}
\begin{algorithmic}[1]
\Require câu dịch $t_1..t_n$, khung $T_1^{\mathrm{orig}}..T_n^{\mathrm{orig}}$, giọng $v$
\State $a_i \gets \textsc{TTS}(t_i, v)$ song song; $R_i \gets |a_i| / T_i^{\mathrm{orig}}$
\State cập nhật tốc độ đọc $\hat{c}_v$ của giọng $v$
\For{$k = 1$ \textbf{to} $M$}
  \State $S \gets \{\, i : R_i > 1{,}3 \text{ và } T_i^{\mathrm{orig}} \ge T_{\min} \,\}$
  \If{$S = \emptyset$} \textbf{break} \EndIf
  \State $\{t_i\}_{i \in S} \gets \textsc{LLM\_RútGọnLô}(\{t_i\}_{i \in S})$ \Comment{một lời gọi}
  \State $a_i \gets \textsc{TTS}(t_i, v)$ song song với $i \in S$; tính lại $R_i$
\EndFor
\For{\textbf{each} $i$ với $R_i > 1{,}0$}
  \State $a_i \gets \textsc{Atempo}(a_i, \min(R_i, 1{,}3))$ \Comment{không gọi lại TTS}
\EndFor
\State \Return $a_1..a_n$ \Comment{khoảng lặng do bước ghép, Mục~\ref{sub:concat}}
\end{algorithmic}
\end{algorithm}

\subsection{Ghép âm thanh}
\label{sub:concat}
Một cách ghép đơn giản là chèn trước mỗi câu $i$ một khoảng lặng cố định $b_i - e_{i-1}$. Tuy nhiên, khi câu $i$ vẫn dài hơn khung sau bước điều chỉnh tốc độ, phần tràn đẩy lùi toàn bộ các câu phía sau và độ trễ cộng dồn theo độ dài bài giảng; thực nghiệm sơ bộ với cách này cho thấy âm thanh của bài dài 15 phút trễ tới 96 giây ở cuối bài. Vì vậy, các câu được đặt theo \emph{mốc tuyệt đối}: gọi $d_i$ là thời lượng âm thanh cuối cùng của câu $i$, thời điểm bắt đầu phát là
\begin{equation}
o_i = \max\bigl(b_i,\; o_{i-1} + d_{i-1}\bigr), \qquad o_0 + d_0 \equiv b_0 .
\label{eq:onset}
\end{equation}
Câu tràn chỉ lấn vào khoảng lặng phía sau, không bao giờ chồng lên câu kế tiếp, và ngay khi gặp một khoảng lặng đủ dài thì các câu sau trở về đúng mốc gốc. Sau khi xếp, nếu âm thanh của đoạn ngắn hơn $L$ thì được đệm lặng; nếu dài hơn (đoạn dày lời, không đủ khoảng lặng để bù) thì cả đoạn được tăng tốc tối đa 15\% bằng bộ lọc \texttt{atempo} của FFmpeg (giữ cao độ). Nhờ đó mỗi đoạn có thời lượng đúng bằng $L$ và độ trễ không lan sang đoạn kế tiếp. Ngoài ra, mốc thời gian do bộ nhận dạng trả về được kẹp trong biên của đoạn, vì thực nghiệm cho thấy mô hình đôi khi trả mốc vượt quá độ dài âm thanh.

% =====================================================================
\section{Triển khai trong hệ thống}
\label{sec:impl}

\subsection{Kiến trúc}
Pipeline được triển khai trong nền tảng LinguaLearn theo kiến trúc hai dịch vụ: dịch vụ nghiệp vụ viết bằng Spring Boot quản lý khóa học, quyền truy cập, hạn ngạch và trạng thái tác vụ; dịch vụ AI viết bằng Python (FastAPI, Celery) thực thi pipeline. Hai dịch vụ giao tiếp qua hàng đợi Redis; tiến độ từng đoạn được phát qua Redis Pub/Sub và chuyển tới trình duyệt bằng WebSocket. Âm thanh được lưu trên dịch vụ lưu trữ đối tượng tương thích S3.

\subsection{Chống trùng và tái sử dụng kết quả}
Do chi phí gọi dịch vụ AI tỷ lệ với thời lượng video, hệ thống áp dụng ba lớp tiết kiệm:
\begin{itemize}
\item \emph{Tái sử dụng bản lồng tiếng}: mỗi cặp (bài giảng, ngôn ngữ đích) chỉ được lồng tiếng một lần; mọi học viên chọn cùng ngôn ngữ dùng chung kết quả và không bị tính vào hạn ngạch.
\item \emph{Khóa chống trùng}: khi nhiều học viên yêu cầu cùng lúc, khóa phân tán (Redis \texttt{SETNX}) theo cặp (bài giảng, ngôn ngữ) bảo đảm chỉ một tác vụ được tạo.
\item \emph{Tái sử dụng bản ký âm}: bản ký âm gốc được dùng lại cho các ngôn ngữ đích khác, bỏ qua bước nhận dạng giọng nói.
\end{itemize}
Ngoài ra, mỗi học viên có hạn ngạch 15 tác vụ lồng tiếng mới mỗi ngày (giảng viên 30) để kiểm soát chi phí. Khi một đoạn phải chạy lại, các câu đã có tệp giọng hợp lệ cho đúng văn bản được dùng lại thay vì gọi TTS lần nữa.

\subsection{Giữ nhạc nền (tùy chọn)}
Một tính năng thử nghiệm tách giọng nói và nhạc nền của video gốc bằng Demucs chạy trên GPU đám mây, rồi trộn nhạc nền với bản lồng tiếng theo tỉ lệ âm lượng gốc, để bài giảng có nhạc mở đầu hay hiệu ứng âm thanh không bị mất phần này. Nhạc nền được chia thành các phần khoảng ba phút, tách song song trên nhiều GPU rồi nối lại bằng trộn chéo. Bước trộn diễn ra sau khi các câu đã được xếp lên trục thời gian, nên không làm thay đổi vị trí hay thời lượng của câu lồng tiếng; tính năng được bật trong toàn bộ thực nghiệm, và thời gian xử lý báo cáo đã gồm bước tách nhạc.

\subsection{Trình phát kép}
Thay vì ghép âm thanh vào video (tốn thời gian mã hóa lại và nhân bản dung lượng theo số ngôn ngữ), hệ thống dùng trình phát kép (Dual Player): video gốc được tắt tiếng và phát song song với tệp âm thanh lồng tiếng. Hai luồng được đồng bộ theo chu kỳ 250 ms; khi độ lệch vượt ngưỡng, vị trí âm thanh được hiệu chỉnh theo video. Cơ chế này hoạt động với cả video tải lên và video YouTube thông qua IFrame Player API. Khi người học tua tới đoạn chưa lồng tiếng xong, trình phát dùng tiếng gốc kèm nhãn báo thay vì để im lặng.

% =====================================================================
\section{Thực nghiệm và đánh giá}
\label{sec:exp}

\subsection{Thiết lập thực nghiệm}
Tập dữ liệu gồm bốn bài giảng tiếng Anh công khai thuộc lĩnh vực công nghệ thông tin (ba bài của loạt Crash Course Computer Science và 15 phút đầu bài giảng Stanford CS229), tổng thời lượng 50 phút, 780 câu, được lồng tiếng sang tiếng Việt. Các tham số của pipeline được liệt kê trong Bảng~\ref{tab:params}. Thực nghiệm gồm hai phần, đều chạy trên một máy chủ ảo (2 vCPU Intel Xeon Skylake, 3,7 GiB RAM, Ubuntu 22.04) dùng chung với hệ thống đang vận hành: (i) tập lệnh đo gọi trực tiếp các hàm của pipeline trên âm thanh của bốn bài giảng để so sánh bốn phương án (mỗi phương án một lần chạy); (ii) bốn tác vụ lồng tiếng thật qua nền tảng, mỗi bài một tác vụ, để đo số liệu vận hành. Bốn phương án so sánh gồm: dịch không ràng buộc độ dài với tốc độ đọc mặc định (V1); dịch có ngân sách thời lượng, tốc độ mặc định (V2); V2 thêm giới hạn số ký tự (\ref{eq:maxchars}) và tăng tốc \texttt{atempo} nhưng không rút gọn (V3); và phương án đề xuất đầy đủ (V4). V3 và V4 dùng chung một bản dịch nên chỉ khác nhau ở bước rút gọn. Toàn bộ thực nghiệm dùng cùng một phiên bản mã nguồn (commit \texttt{0d866b9}), mô hình Gemini 3.6 Flash với mức suy nghĩ \texttt{minimal}, và bật tính năng giữ nhạc nền (Mục~\ref{sec:impl}). Tốc độ đọc $\hat{c}_v$ của giọng được xóa trước khi đo, nên hai bài chạy đầu tiên được dịch khi giọng chưa có số đo (không có giới hạn ký tự), hai bài còn lại dùng số đo từ bài trước; điều này phản ánh đúng hoạt động thật của hệ thống.

\begin{table}[t]
\caption{Tham số của pipeline}
\label{tab:params}
\centering
\setlength{\tabcolsep}{4pt}
\begin{tabular}{|l|l|}
\hline
\textbf{Tham số} & \textbf{Giá trị} \\
\hline
Độ dài đoạn $L$ & 600 giây \\
Số lần thử lại mỗi đoạn & 3 \\
Ngưỡng tỷ lệ lời thoại tối thiểu & 10\% \\
Ngưỡng $R_{\min}$ / $R_{\max}$ & 1,0 / 1,3 \\
Tốc độ đọc tăng tối đa & +30\% \\
Số lần rút gọn tối đa $M$ & 2 \\
Khung tối thiểu để rút gọn $T_{\min}$ & 2 giây \\
Tăng tốc cả đoạn tối đa & $\times$1,15 \\
Mô hình nhận dạng giọng nói & Whisper large-v3-turbo \\
Mô hình ngôn ngữ lớn & Gemini 3.6 Flash \\
Mức suy nghĩ của LLM (dịch, rút gọn) & \texttt{minimal} \\
Hệ số giới hạn ký tự $\rho$ & 1,15 \\
Tổng hợp giọng nói & Azure Speech \\
Tăng tốc câu & FFmpeg \texttt{atempo} \\
Số yêu cầu TTS đồng thời & 16 \\
\hline
\end{tabular}
\end{table}

\subsection{Độ đo}
\begin{itemize}
\item \emph{Phân bố tỷ số $R$}: tỷ lệ câu rơi vào từng nhánh của (\ref{eq:policy}) khi dịch không ràng buộc và khi dịch có ràng buộc độ dài.
\item \emph{Độ lệch thời lượng}: $\Delta_i = |T_i^{\mathrm{final}} - T_i^{\mathrm{orig}}|$ trung bình theo câu, so sánh giữa phương án phát nguyên tốc độ và phương án đề xuất.
\item \emph{Độ trễ theo trục thời gian}: sau khi ghép cả bài, độ lệch $o_i - b_i$ giữa thời điểm phát của câu lồng tiếng và mốc gốc (trung bình, p95, cực đại) và chênh lệch tổng thời lượng so với video gốc. Độ đo này bắt được trễ cộng dồn mà $\Delta_i$ theo từng câu không phản ánh, và được dùng để so sánh cách ghép khoảng lặng cố định với cách xếp theo mốc tuyệt đối (\ref{eq:onset}).
\item \emph{Thời gian xử lý}: thời gian thực thi trên mỗi phút video và thời gian đến khi có đoạn đầu tiên.
\item \emph{Chi phí}: số lần gọi LLM và TTS.
\item \emph{Độ bảo toàn ý của câu rút gọn}: 50 câu bị rút gọn được chọn ngẫu nhiên (hạt giống cố định) và chấm trên thang 0--2 (2: giữ đủ ý, 1: mất ý phụ, 0: mất hoặc sai ý chính) bởi hai thành viên nhóm nghiên cứu, chấm độc lập; sau đó hai người thống nhất điểm cho các câu chấm lệch, dùng bản dịch đầy đủ trước khi rút gọn làm chuẩn so sánh. Báo cáo độ đồng thuận ban đầu và điểm sau khi thống nhất.
\end{itemize}

\subsection{Kết quả}
\label{sub:results}

\paragraph{Độ khớp thời lượng của từng câu.} Bảng~\ref{tab:results} so sánh bốn phương án trên 780 câu. Đưa ngân sách thời lượng vào lời nhắc dịch giảm độ lệch trung bình từ 1,89 xuống 1,25 giây; thêm giới hạn số ký tự và tăng tốc \texttt{atempo} trong giới hạn 30\% giảm tiếp xuống 0,29 giây; thêm bước rút gọn theo lô (phương án đề xuất) còn 0,19 giây, và chỉ 30,4\% số câu còn dài hơn khung gốc so với 96,9\% khi dịch không ràng buộc. Phương án đề xuất chỉ cần 20 lời gọi LLM (7 lời gọi dịch, mỗi đoạn một lời gọi, và 13 lời gọi rút gọn theo lô) và 1.034 lần tổng hợp giọng nói cho 780 câu. Khi chạy lại phương án đề xuất lần thứ hai, lúc giọng đọc đã có số đo tốc độ, độ lệch trung bình là 0,194 giây, với 21 lời gọi LLM và 845 lần tổng hợp giọng nói.

\begin{table}[t]
\caption{So sánh độ khớp thời lượng từng câu giữa các phương án (780 câu, mỗi phương án một lần chạy)}
\label{tab:results}
\centering
\setlength{\tabcolsep}{3pt}
\begin{tabular}{|l|c|c|c|c|c|}
\hline
\textbf{PA} & \textbf{$\bar{\Delta}$ (s)} & \textbf{Tràn (\%)} & \textbf{$R>1{,}3$ (\%)} & \textbf{LLM} & \textbf{TTS} \\
\hline
V1 & 1,89 & 96,9 & 82,8 & 7 & 780 \\
V2 & 1,25 & 90,5 & 63,7 & 7 & 780 \\
V3 & 0,29 & 46,5 & 23,2 & 7 & 780 \\
V4 & 0,19 & 30,4 & 19,6 & 20 & 1.034 \\
\hline
\end{tabular}
\end{table}

\paragraph{Độ trễ theo trục thời gian.} Bảng~\ref{tab:timeline} cho độ lệch tuyệt đối $|o_i - b_i|$ giữa thời điểm phát của câu lồng tiếng và mốc gốc, đo trên âm thanh được ghép lại từ các tệp câu thật bằng đúng hàm xếp câu của pipeline. Ở cả hai phương án, tổng thời lượng sau khi ghép bằng video gốc (chênh không quá 0,1 giây), vì mỗi đoạn được khớp đúng độ dài $L$. Tuy nhiên, khi không điều chỉnh tốc độ (V2), phần tràn vượt xa khoảng lặng có sẵn: độ trễ trung bình là 14,9 giây, câu muộn nhất trễ 58,1 giây, và mỗi đoạn phải tăng tốc tối đa $\times 1{,}15$ rồi vẫn bị cắt tới 93,4 giây phần cuối. Với phương án đề xuất, độ trễ trung bình chỉ còn 0,40 giây, phân vị 95 là 1,98 giây, 10,8\% số câu lệch quá 1 giây, hệ số tăng tốc cả đoạn lớn nhất là 1,002 và phần cuối bị cắt không quá 0,01 giây. Bài giảng trên lớp (15 phút, tỷ lệ lời thoại thấp nhất và nhiều câu ngắn) vẫn là bài khó nhất với độ trễ trung bình 0,68 giây và 21,6\% số câu lệch quá 1 giây; ba bài còn lại có độ trễ trung bình từ 0,11 đến 0,30 giây.

\begin{table}[t]
\caption{Độ trễ theo trục thời gian trên bốn bài (780 câu), xếp câu theo mốc tuyệt đối. Cột TB, P95 là $|o_i - b_i|$; cột Max là độ trễ muộn nhất, tính bằng giây}
\label{tab:timeline}
\centering
\setlength{\tabcolsep}{3pt}
\begin{tabular}{|l|c|c|c|c|c|}
\hline
\textbf{PA} & \textbf{Chênh (s)} & \textbf{TB} & \textbf{P95} & \textbf{Max} & \textbf{$>$1 s (\%)} \\
\hline
V2 & $-0{,}02\ldots0{,}08$ & 14,88 & 44,24 & 58,08 & 93,1 \\
V4 & $-0{,}02\ldots0{,}08$ & 0,40 & 1,98 & 6,41 & 10,8 \\
\hline
\end{tabular}
\end{table}

\paragraph{Vận hành trên hệ thống thật.} Bốn bài được lồng tiếng lại qua giao diện web sau khi xóa toàn bộ dữ liệu đã lưu của chúng (bản ký âm, bản dịch, bản lồng tiếng, nhạc nền đã tách), nên đây là lần lồng tiếng đầu tiên, có đủ bước nhận dạng giọng nói và tách nhạc. Cả bốn tác vụ hoàn tất, không có lỗi hay lần thử lại nào, và thời lượng âm thanh bằng đúng thời lượng bài học (Bảng~\ref{tab:server}). Thời gian xử lý trung bình là 16,2 giây cho mỗi phút video (từ 14,5 đến 19,3), tức nhanh hơn thời gian thực khoảng 3,7 lần, trên máy chủ 2 vCPU dùng chung với hệ thống đang vận hành. Người học nghe được đoạn đầu tiên sau 148 đến 181 giây. Ở ba bài có hai đoạn, nhờ xử lý theo dây chuyền, đoạn thứ hai hoàn tất chỉ 25 đến 50 giây sau đoạn đầu. So với phiên bản trước của hệ thống trên cùng bốn bài và cùng giọng đọc (dịch hai lời gọi, tăng tốc bằng cách gọi lại TTS, rút gọn từng câu, chưa có dây chuyền), thời gian xử lý giảm từ khoảng 110 giây xuống 16,2 giây cho mỗi phút video và thời gian tới đoạn đầu giảm từ khoảng 1.045 giây xuống 168 giây. Trong 760 câu của bốn tác vụ, 27,4\% câu vẫn còn $R > 1{,}3$ khi tăng tốc, cao nhất ở bài giảng trên lớp (45\%).

\begin{table}[t]
\caption{Bốn tác vụ lồng tiếng thật trên nền tảng (lần lồng tiếng đầu, bật giữ nhạc nền)}
\label{tab:server}
\centering
\setlength{\tabcolsep}{3pt}
\begin{tabular}{|l|c|c|c|c|c|}
\hline
\textbf{Bài} & \textbf{Dài (s)} & \textbf{Đoạn} & \textbf{Xong (s)} & \textbf{s/phút} & \textbf{Đoạn đầu (s)} \\
\hline
Bài 1 & 566 & 1 & 181,8 & 19,3 & 181,4 \\
Bài 2 & 815 & 2 & 223,7 & 16,5 & 173,9 \\
Bài 3 & 718 & 2 & 173,1 & 14,5 & 148,0 \\
Bài 4 & 900 & 2 & 220,1 & 14,7 & 170,3 \\
\hline
\end{tabular}
\end{table}

\paragraph{Độ bảo toàn ý của câu rút gọn.} Trong 165 câu được rút gọn của phương án đề xuất, 50 câu được chọn ngẫu nhiên để chấm. Hai người chấm độc lập cho tỷ lệ giữ đủ ý lần lượt là 58\% và 74\%, mất ý phụ 38\% và 26\%, mất hoặc sai ý chính 4\% và 0\%. Hai người cho cùng điểm ở 26/50 câu (52\%; AC1 của Gwet 0,38; kappa của Cohen xấp xỉ 0 vì phần lớn câu được chấm điểm 2). Khác biệt chủ yếu đến từ cách hiểu ``ý phụ'': một người trừ điểm khi mất bất kỳ chi tiết có nội dung nào, người kia trừ điểm cả khi chỉ bỏ từ đệm. Sau khi thống nhất 24 câu lệch theo quy tắc chung (2: chỉ bỏ từ đệm hoặc thay từ đồng nghĩa; 1: mất một chi tiết có nội dung nhưng ý chính còn đúng; 0: mất hoặc đổi ý chính, hoặc câu bị cụt đến mức khó hiểu), 66\% câu giữ đủ ý, 30\% mất ý phụ và 4\% (2 câu) mất ý chính, điểm bảo toàn ý trung bình 1,62 trên 2. Hai câu điểm 0 đều là câu bị rút gọn đến mức cụt, mất một vế liệt kê hoặc động từ chính; các lỗi điểm 1 thường gặp là mất thuật ngữ (``học máy''), mất đối tượng của hành động, hoặc thay thuật ngữ bằng từ kém chuẩn (``chương trình'' thành ``bài'').

\subsection{Thảo luận}
Mỗi thành phần của phương pháp đóng góp vào độ khớp thời lượng: ngân sách thời lượng trong lời nhắc giảm độ lệch trung bình khoảng 34\% so với dịch không ràng buộc (1,89 xuống 1,25 giây), giới hạn số ký tự cùng tăng tốc \texttt{atempo} giảm thêm khoảng 77\% (còn 0,29 giây), và rút gọn theo lô giảm thêm khoảng 36\% (còn 0,19 giây). Cách xếp theo mốc tuyệt đối là điều kiện cần nhưng không đủ: không có điều chỉnh tốc độ, phần tràn vượt quá khoảng lặng có sẵn nên vẫn có câu trễ gần một phút và phần cuối mỗi đoạn bị cắt; chỉ khi kết hợp với giới hạn độ dài, tăng tốc và rút gọn, độ trễ mới nằm dưới một giây ở phần lớn các câu.

Gộp các bước gọi mô hình là thay đổi chính về chi phí. Mỗi đoạn chỉ cần một lời gọi dịch, các câu cần rút gọn được gửi chung trong một lời gọi cho mỗi vòng, và tăng tốc được thực hiện trên tệp có sẵn thay vì gọi lại TTS. Với cùng bốn bài, phiên bản trước của hệ thống cần 286 lời gọi LLM và 1.678 lần tổng hợp giọng nói, phiên bản hiện tại cần 20 và 1.034. Giới hạn số ký tự ước theo tốc độ đọc đo được giúp bản dịch vừa khung ngay từ đầu, thể hiện ở việc tỷ lệ câu có $R > 1{,}3$ giảm từ 63,7\% (V2) xuống 48,1\% trước khi tăng tốc. Quy tắc không rút gọn câu có khung dưới 2 giây là một đánh đổi có chủ đích: các câu này thường là lời chào hay câu cảm ơn ngắn, phần lớn rút gọn vẫn vượt ngưỡng, và phần dư được bù bằng khoảng lặng phía sau.

Các trường hợp còn lỗi gồm: bài giảng trên lớp có tỷ lệ lời thoại thấp, nhiều câu ngắn, nên còn nhiều câu vượt ngưỡng và độ trễ lớn hơn hẳn; câu bị rút gọn đôi khi mất thuật ngữ hoặc bị cụt (Mục độ bảo toàn ý); và mốc thời gian của bộ nhận dạng giọng nói đôi khi vượt quá độ dài âm thanh, đã được xử lý bằng cách kẹp mốc.

Các hạn chế của nghiên cứu gồm: (i) chưa đánh giá cảm nhận của người nghe về độ tự nhiên, độ dễ hiểu và độ khớp hình, trong khi tăng tốc bằng \texttt{atempo} là kéo giãn tín hiệu, nên kết luận của bài chỉ về độ khớp thời lượng và độ bảo toàn ý, không về chất lượng nghe; (ii) tập dữ liệu nhỏ (bốn bài, 50 phút, tiếng Anh, lĩnh vực công nghệ thông tin) và mỗi phương án chỉ chạy một lần, trong khi LLM không tất định; (iii) thời gian xử lý đo trên một máy chủ 2 vCPU dùng chung với hệ thống đang vận hành, mỗi bài một tác vụ, và phụ thuộc vào thời gian chờ cấp GPU đám mây cho bước tách nhạc; (iv) hiệu quả tái sử dụng bản ký âm và bản lồng tiếng giữa các ngôn ngữ đích chưa được đo; (v) độ bảo toàn ý do hai thành viên nhóm chấm trên 50 câu, độ đồng thuận ban đầu thấp và điểm cuối có được nhờ thống nhất.

% =====================================================================
\section{Kết luận}
\label{sec:conclusion}
Bài báo đã trình bày một pipeline lồng tiếng tự động cho video bài giảng được tích hợp trong hệ thống quản lý học tập. Bằng cách xử lý ràng buộc thời lượng ở cả tầng prompt (ngân sách thời lượng và giới hạn số ký tự ước theo tốc độ đọc đo được) và tầng hậu xử lý (rút gọn theo lô và tăng tốc tín hiệu có giới hạn), hệ thống tạo ra âm thanh lồng tiếng có thời lượng khớp với video gốc mà không cần huấn luyện mô hình riêng: độ lệch thời lượng trung bình của mỗi câu giảm từ 1,89 xuống 0,19 giây và độ trễ tuyệt đối trung bình sau khi ghép cả bài là 0,40 giây, với 20 lời gọi LLM cho 50 phút bài giảng. Kiến trúc phân đoạn mười phút xử lý theo dây chuyền cùng các cơ chế chống trùng và tái sử dụng giúp hệ thống chịu lỗi cục bộ và phản hồi sớm: trên hệ thống thật, bốn tác vụ lồng tiếng hoàn tất không lỗi, mỗi phút video được xử lý trong khoảng 16 giây và người học nghe được đoạn đầu sau khoảng ba phút; hiệu quả tái sử dụng chưa được đo.

Hạn chế hiện tại là giọng đọc là giọng tổng hợp có sẵn, chưa giữ được âm sắc của giảng viên; đánh giá mới thực hiện trên tập dữ liệu nhỏ, mỗi phương án một lần chạy, và chưa có đánh giá cảm nhận của người nghe. Hướng phát triển gồm: khảo sát điểm đánh giá trung bình (MOS) với nhóm người nghe độc lập về độ tự nhiên và độ khớp hình, nhân bản giọng nói của giảng viên, đồng bộ khẩu hình, mở rộng sang các lĩnh vực khác như kỹ năng mềm, và so sánh hiệu quả học tập giữa học bằng phụ đề và học bằng lồng tiếng.

\section*{Lời cảm ơn}
Nhóm tác giả xin cảm ơn Khoa Công nghệ Thông tin, Trường Đại học Công nghiệp Thành phố Hồ Chí Minh đã tạo điều kiện thực hiện nghiên cứu này.

% =====================================================================
\begin{thebibliography}{00}
\bibitem{mayer2009} R. E. Mayer, \textit{Multimedia Learning}, 2nd ed. Cambridge, U.K.: Cambridge Univ. Press, 2009.
\bibitem{kalyuga1999} S. Kalyuga, P. Chandler, and J. Sweller, ``Managing split-attention and redundancy in multimedia instruction,'' \textit{Appl. Cogn. Psychol.}, vol. 13, no. 4, pp. 351--371, 1999.
\bibitem{radford2023whisper} A. Radford, J. W. Kim, T. Xu, G. Brockman, C. McLeavey, and I. Sutskever, ``Robust speech recognition via large-scale weak supervision,'' in \textit{Proc. 40th Int. Conf. Mach. Learn. (ICML)}, Honolulu, HI, USA, 2023, pp. 28492--28518.
\bibitem{gemini2023} Gemini Team, Google, ``Gemini: A family of highly capable multimodal models,'' arXiv:2312.11805, 2023.
\bibitem{federico2020} M. Federico, R. Enyedi, R. Barra-Chicote, R. Giri, U. Isik, A. Krishnaswamy, and H. Sawaf, ``From speech-to-speech translation to automatic dubbing,'' in \textit{Proc. 17th Int. Conf. Spoken Lang. Transl. (IWSLT)}, 2020, pp. 257--264.
\bibitem{brannon2023} W. Brannon, Y. Virkar, and B. Thompson, ``Dubbing in practice: A large scale study of human localization with insights for automatic dubbing,'' \textit{Trans. Assoc. Comput. Linguist.}, vol. 11, pp. 419--435, 2023.
\bibitem{oktem2019} A. \"{O}ktem, M. Farr\'{u}s, and A. Bonafonte, ``Prosodic phrase alignment for machine dubbing,'' in \textit{Proc. Interspeech}, Graz, Austria, 2019, pp. 4215--4219.
\bibitem{lakew2022} S. M. Lakew, Y. Virkar, P. Mathur, and M. Federico, ``Isometric MT: Neural machine translation for automatic dubbing,'' in \textit{Proc. IEEE Int. Conf. Acoust., Speech Signal Process. (ICASSP)}, Singapore, 2022, pp. 6242--6246.
\bibitem{wu2023videodubber} Y. Wu, J. Guo, X. Tan, C. Zhang, B. Li, R. Song, L. He, S. Zhao, A. Menezes, and J. Bian, ``VideoDubber: Machine translation with speech-aware length control for video dubbing,'' in \textit{Proc. AAAI Conf. Artif. Intell.}, vol. 37, no. 11, 2023, pp. 13772--13779.
\bibitem{bain2023whisperx} M. Bain, J. Huh, T. Han, and A. Zisserman, ``WhisperX: Time-accurate speech transcription of long-form audio,'' in \textit{Proc. Interspeech}, Dublin, Ireland, 2023, pp. 4489--4493.
\end{thebibliography}

\end{document}
